\section{Introduction}
\label{sec:introduction}

Mass spectrometry imaging (MSI) records a full mass spectrum at each location
of a tissue section, mapping many molecules at once without choosing targets in
advance~\citep{caprioli1997maldi,schwamborn2010molecular}. The data are large:
the sections used here have 972--4{,}761 measured pixels and 1{,}481--85{,}062
mass-to-charge ($m/z$) bins. Most analyses therefore begin by selecting a small
set of informative peaks. \emph{Peak picking} applies rule-based thresholds to
each spectrum and can discard weak but informative
ions~\citep{gibb2012maldiquant,lieb2020sparse}. \emph{Peak learning} instead
fits an unsupervised model and ranks $m/z$ bins by what the model uses.
Examples are msiPL, a variational autoencoder (VAE) on single
spectra~\citep{abdelmoula2021peak}, and S3PL, which adds spatial context through
3D convolution over tissue patches~\citep{weigand2026spatial}.

The proposal for this project, \emph{Representation Learning for Autonomous
Peak Detection in MSI}, set out to build \emph{Spatial-msiPL}: an msiPL-style
VAE given cheap local spatial context (the mean of a $3\times3$ neighbourhood),
as a lightweight alternative to 3D convolution. Testing it fairly raised a
second question: once a model is fitted, which part of the pipeline determines
peak quality? The candidates are the representation, the segmentation built on
it, and the attribution method used to rank bins.

We built Spatial-msiPL and its variants (uniform and attention-weighted
neighbourhoods, larger windows, masked centres) and segmented each latent space
with a Gaussian mixture model (GMM). Peaks were ranked by Integrated Gradients
(IG)~\citep{sundararajan2017axiomatic} of the GMM posteriors, or by directly
correlating each ion image with the posterior maps. All methods were scored at
matched peak budgets under rules fixed before the results, on 24 sections from
three public collections. The comparison included reproduced msiPL and S3PL
baselines and model-free rankings.

Holding the segmentation fixed, correlation outperformed IG in 136 of 140
paired comparisons across three training seeds, and this was statistically
supported in all six dataset--model comparisons. Neighbourhood input to the
encoder gave no consistent, spatially specific gain. Spatially aware
\emph{ranking} statistics, namely segmentation correlation and Moran's~$I$,
were the strongest label-free rankings. A linear PCA segmentation matched the
VAE. S3PL's published glioblastoma score lay outside ten reproduced seeds, and
weight initialisation switched individual S3PL sections between distinct
solutions.

The contributions are:
(i) a predeclared, matched-budget evaluation of learned MSI peak ranking on
three datasets;
(ii) a ranking ablation, with paired statistics, showing that
segment-then-correlate outperforms gradient attribution in this pipeline;
(iii) controlled tests that bound where spatial information helps;
(iv) a statistical account of S3PL's reproduction gap and seed sensitivity.
\Cref{sec:background} reviews related work, \cref{sec:questions} states the
research questions, \cref{sec:methodology} gives the methods,
\cref{sec:results} the results, \cref{sec:discussion} the discussion, and
\cref{sec:conclusion} concludes.
